\chapter{État de l'art et fondements métiers}

\section*{Introduction}
\plannedcontent{Présenter le rôle de ce chapitre : situer le stage, introduire la gestion
des infrastructures de recharge, formuler la problématique et justifier la solution.}

\section{Cadre du projet}

\subsection{Cadre académique}
\plannedcontent{Stage d'été réalisé après la première année du cycle ingénieur en génie
logiciel (IGL3), au Département des Sciences de l'Informatique de la Faculté des Sciences
de Tunis. Présenter les objectifs pédagogiques et les compétences mobilisées.}

\subsection{Présentation de l'organisme d'accueil}
\plannedcontent{Présenter \CompanyName{}, son domaine d'activité, son organisation et le
contexte dans lequel le sujet a été proposé. Les informations institutionnelles seront
validées avec M. \CompanySupervisor{} avant rédaction définitive.}

\subsection{Présentation du domaine d'application}

\subsubsection{Mobilité électrique et infrastructures de recharge}
\plannedcontent{Introduire les enjeux de la mobilité électrique, les réseaux de recharge,
les bornes, les connecteurs et les principales parties prenantes.}

\subsubsection{Supervision d'un réseau de bornes}
\plannedcontent{Expliquer le rôle d'un CSMS, la remontée d'événements, le suivi des
sessions, la maintenance et la nécessité d'une information opérationnelle fiable.}

\subsubsection{OCPP et interopérabilité}
\plannedcontent{Présenter OCPP 1.6 JSON, les échanges WebSocket, les messages de statut,
les heartbeats et les commandes distantes utiles au projet.}

\subsubsection{Disponibilité et qualité de service}
\plannedcontent{Définir les notions de disponible, occupée, en défaut, hors ligne et en
maintenance. Introduire les règles de projection et les délais de connectivité.}

\section{Problématique}
\plannedcontent{Décrire la difficulté de connaître l'état réel d'une borne, la fragmentation
des informations, les pannes silencieuses, les effets sur les conducteurs et la nécessité
d'une coordination multi-rôle et multi-organisation.}

\section{Étude de l'existant}

\subsection{Solutions existantes}
\plannedcontent{Étudier des plateformes de gestion et de recharge comparables, leurs
parcours de recharge, leur supervision, leur tarification et leurs mécanismes de
maintenance. Les solutions et critères seront documentés avec des sources à jour.}

\subsection{Standards, outils et simulateurs}
\plannedcontent{Comparer les outils de simulation OCPP, les solutions cartographiques,
les approches de paiement et les moyens de notification pertinents pour un MVP.}

\subsection{Synthèse comparative}
\plannedcontent{Présenter un tableau de comparaison : temps réel, OCPP, cartographie,
multi-tenant, recharge, paiement, alertes, maintenance, rapports et extensibilité.}

\section{Solution proposée}

\subsection{Espace public et création d'accès}
\plannedcontent{Landing page, demande de démonstration, inscription client, OAuth Google,
invitations et activation des comptes administrateur, opérateur et technicien.}

\subsection{Espaces authentifiés par rôle}
\plannedcontent{Présenter les responsabilités du Super Administrateur, de l'Administrateur,
de l'Opérateur, du Technicien et du Client.}

\subsection{Supervision OCPP et simulation}
\plannedcontent{Expliquer le rôle du gateway OCPP, du simulateur de borne, des projections
de disponibilité et du simulateur de paiement dans l'environnement de développement.}

\section{Méthodologie de travail}

\subsection{Approche itérative et incrémentale}
\plannedcontent{Décrire une démarche individuelle inspirée de l'agilité : backlog,
incréments verticaux, démonstration, validation et amélioration continue, sans prétendre
appliquer toutes les cérémonies d'une équipe Scrum complète.}

\subsection{Organisation et suivi}
\plannedcontent{Présenter le découpage en blocs, la gestion Git, les validations
fonctionnelles, les tests automatisés et le suivi de l'avancement.}

\section*{Conclusion}
\plannedcontent{Synthétiser le contexte et introduire la spécification détaillée des
besoins.}

